\section{Character Sums via the Quadratic Character}\label{sec:quadratic-char}
For certain special characters, the associated Gaussian sums can be evaluated explicitly. A very special multiplicative character to study is the quadratic character, $\eta$. We get a very celebrated formula for the Gaussian sum of $\eta$ and the canonical additive character. 
\begin{Theorem}{Gaussian Sum of the Quadratic Character}{}
  \index{quadratic character!Gaussian sum evaluation}%
  Let $\bbF_q$ be a finite field with $q=p^r$ where $p$ is an odd prime and $r\in\bbN$. Let $\eta$ be the quadratic character of $\bbF_q$ and let $\chi_1$ be the canonical additive character of $\bbF_q$. Then $$G(\eta,\chi_1)=\begin{cases}
    (-1)^{r-1}\cdot q^{\nicefrac12} & \text{if }p\equiv 1\bmod 4\\ 
    (-1)^{r-1}\cdot i^{r}\cdot q^{\nicefrac12} & \text{if }p\equiv 3\bmod 4    
  \end{cases}$$
  In other words, $$G(\eta,\chi_1)=(-1)^{r-1}\cdot i^{\frac{(p-1)^2}4 r} \cdot q^{\nicefrac12}$$ 
\end{Theorem}
\begin{proof}
  Since $\eta$ always takes values in $\pm 1$ we have $\eta=\ov{\eta}$. Hence by using \lemref[(iv)]{lem:gaussian-sum-identity} we have $G(\eta,\chi_1)^2=\eta(-1)q$. Now if $q\equiv 1\bmod 4$ then $\eta(-1)=1$ and if $q\equiv 3\bmod 4$ then $\eta(-1)=-1$. Hence \begin{equation}
    G(\eta,\chi_1)=\begin{cases}
      \pm\ q^{\nicefrac12}& \text{if }q\equiv 1\bmod 4\\
      \pm\ i\cdot q^{\nicefrac12}& \text{if }q\equiv 3\bmod 4
    \end{cases}
  \end{equation}
  Since $i^{\nicefrac{(p-1)^2}{4}}=\begin{cases}
      1 & \text{if }q\equiv 1\bmod 4\\ 0 & \text{if }q\equiv 3\bmod 4
    \end{cases}$

  \begin{observation}\label{obs:gauss-quad-char}
    For any finite field $\bbF_q$,  $G(\eta,\chi_1)=\pm\ i^{\nicefrac{(q-1)^2}{4}}\cdot q^{\nicefrac12}$
  \end{observation}

  It remains to resolve the $\pm$ sign. We first prove this for $r=1$ then extend to $r>1$.

  \paragraph*{Step 1: Setting up the linear operator $T$.}
  Let $W$ be the vector space of all complex valued functions on $\bbF_p^*$, a $(p-1)$-dimensional vector space over $\bbC$. Let $\mcB\coloneqq \{f_1,\dots, f_{p-1}\}$ where $f_j(k)=1$ if $j=k$ and otherwise $0$ be the standard basis of $W$. By the orthogonality of multiplicative characters in \thmref{thm:mult-char-props} the set $\mcB_{\Mq}\coloneqq \{\psi_0,\dots, \psi_{p-2}\}$ as defined in \lemref{lem:mult-char-define} also forms a basis of $W$.

  Let $\zeta_p=e^{\nicefrac{2\pi i}{p}}$ be the $p^{th}$ root of unity. Hence if $\chi_k\in\Xp$ is as defined in \thmref{thm:add-char-define}, then $\chi_k(j)=\zeta_p^{jk}$ for all $j\in\bbF_p$. Define the linear operator $T:W\to W$ by
  \begin{equation}\label{eq:quad-linear-operator}
    \forall\ h\in W,\ \forall\ k\in[p-1]\qquad (Th)(k)=\sum_{j=1}^{p-1}\zeta_p^{jk}h(j).
  \end{equation}
  For each basis element $\psi_i$, $$(T\psi_i)(k)=\sum_{j=1}^{p-1}\chi_k(j)\cdot \psi_i(j)=G(\psi_i,\chi_k)=\ov{\psi_i}(k)\cdot G(\psi_i,\chi_1)\quad \forall\ k\in\bbF_p^*,$$where the last equality uses \lemref[(i)]{lem:gaussian-sum-identity}. Hence $T\psi_i=G(\psi_i,\chi_1)\cdot \ov{\psi_i}$ for every multiplicative character of $\bbF_p$.

  \paragraph*{Step 2: Computing $\det(T)$ via the character basis.}
  Since $\ov{\psi}=\psi$ only for $\psi_0$ and $\eta=\psi_{\nicefrac{p-1}2}$, the matrix of $T$ in $\mcB_{\Mq}$ has two diagonal entries $G(\psi_0,\chi_1)=-1$ and $G(\eta,\chi_1)$; all other characters pair up as conjugates $(\psi,\ov{\psi})$, each contributing the block $\begin{pmatrix}0&G(\ov{\psi},\chi_1)\\G(\psi,\chi_1)&0\end{pmatrix}$ with determinant contribution $-\psi(-1)p$ by \lemref[(iv)]{lem:gaussian-sum-identity}. Using $\psi_j(-1)=(-1)^j$ (since $\psi(-1)=-1$ for the principal character),
  \begin{align*}
    \det(T) & = - G(\eta,\chi_1) \cdot (-p)^{\nicefrac{(p-3)}2}\prod_{j=1}^{\frac{p-3}2}(-1)^j\\
    & =  (-1)^{\nicefrac{(p-1)}{2}}\cdot G(\eta,\chi_1) \cdot p^{\nicefrac{(p-3)}2}\cdot (-1)^{\nicefrac{(p-1)(p-3)}8}\\
    & = \pm (-1)^{\nicefrac{(p-1)}{2}}\cdot i^{\nicefrac{(p-1)(p-2)}2}\cdot p^{\nicefrac{(p-2)}{2}},
  \end{align*}
  where in the last step we substitute \obsref{obs:gauss-quad-char} and simplify the exponent. Thus:
  \begin{equation}\label{eq:fourier-basis-det}
    \det(T)=\pm (-1)^{\nicefrac{(p-1)}{2}}\cdot i^{\nicefrac{(p-1)(p-2)}2}\cdot p^{\nicefrac{(p-2)}{2}}.
  \end{equation}

  \paragraph*{Step 3: Computing $\det(T)$ via the standard basis.}
  From \eqref{eq:quad-linear-operator}, the $(j,k)$ entry of $T$ in $\mcB$ is $\zeta_p^{jk}$. Factoring and reducing modulo $p$ (since $p\mid 1+2+\cdots+(p-1)$), the determinant equals that of the Vandermonde matrix $V(\zeta_p,\zeta_p^2,\dots,\zeta_p^{p-1})$:
  $$\det(T)=\prod_{1\leq j<k\leq p-1}\lt(\zeta_p^k-\zeta_p^j\rt).$$
  Setting $\zeta_{2p}=e^{\nicefrac{\pi i}{p}}$ so that $\zeta_{2p}^2=\zeta_p$ and using $\zeta_{2p}^t-\zeta_{2p}^{-t}=2i\sin\frac{t\pi}{p}$, we expand the Vandermonde product:
  \begin{align*}
    \det(T) & = \prod_{1\leq j<k\leq p-1}\zeta_{2p}^{j+k}\lt(\zeta_{2p}^{k-j}-\zeta_{2p}^{-(k-j)}\rt)\\
    & = \zeta_{2p}^{\sum_{j<k}(j+k)} \cdot i^{\nicefrac{(p-1)(p-2)}{2}}\prod_{1\leq j<k\leq p-1} \lt(2\sin\frac{\pi(k-j)}{p}\rt).
  \end{align*}
  Computing the exponent sum $\sum_{1\leq j<k\leq p-1}(j+k)=\frac{p(p-1)(p-2)}{2}$ gives $\zeta_{2p}^{\sum(j+k)}=(-1)^{\nicefrac{(p-1)}{2}}$, so:
  \begin{equation}\label{eq:standard-basis-det}
    \det(T)=(-1)^{\nicefrac{(p-1)}{2}}\cdot i^{\nicefrac{(p-1)(p-2)}{2}}\cdot \prod_{1\leq j<k\leq p-1} \lt(2\sin\frac{\pi(k-j)}{p}\rt).
  \end{equation}

  \paragraph*{Step 4: Resolving the sign for $r=1$.}
  Comparing \eqref{eq:fourier-basis-det} and \eqref{eq:standard-basis-det}, the first two factors coincide. Since $k-j\in[p-1]$ for all $1\leq j<k\leq p-1$, we have $\sin\frac{\pi(k-j)}{p}>0$, so $\prod_{j<k}2\sin\frac{\pi(k-j)}{p}>0$. Therefore the product on the right of \eqref{eq:standard-basis-det} equals $p^{\nicefrac{(p-2)}{2}}$, and matching with \eqref{eq:fourier-basis-det} forces the $\pm$ to be $+$:
  $$G(\eta,\chi_1)=i^{\nicefrac{(p-1)^2}{4}}\cdot p^{\nicefrac12}.$$

  \paragraph*{Step 5: Extension to $r>1$ via character lifting.}
  For $r>1$, lift $\eta$ and $\chi_1$ from $\bbF_p$ to $\bbF_q$ ($q=p^r$) via $\eta^{(r)}=\eta\circ N$ and $\chi_1^{(r)}=\chi_1\circ\tr$, where $N$ and $\tr$ are the absolute norm and trace from $\bbF_q$ to $\bbF_p$ respectively. By \DHthm,
  $$G(\eta^{(r)},\chi_1^{(r)})=(-1)^{r-1}G(\eta,\chi_1)^r=(-1)^{r-1}\lt(i^{\nicefrac{(p-1)^2}{4}}\cdot p^{\nicefrac12}\rt)^r=(-1)^r\cdot i^{\frac{(p-1)^2r}{4}}\cdot q^{\nicefrac12}.$$
  Since $\eta^{(r)}$ is the quadratic character of $\bbF_q$, this is the stated formula.
\end{proof}

The proof of the above theorem gives us another important result. 
\begin{lemma}{Vandermonde Determinant with Roots of Unity}{}
  If $p$ is a odd prime and $\zeta_p$ is the $p^{th}$ root of unity. Let $V(\zeta_p,\zeta_p^2,\dots,\zeta_p^{p-1})$ is the $p\times p$ Vandermonde Matrix. Then $$\det\lt(V(\zeta_p,\zeta_p^2,\dots,\zeta_p^{p-1})\rt)=(-1)^{\nicefrac{(p-1)}{2}}\cdot i^{\nicefrac{(p-1)(p-2)}2}\cdot p^{\nicefrac{(p-2)}{2}}$$
\end{lemma}

As a striking application of the machinery developed around the quadratic character, we can now give a proof of one of the classical result of  number theory -- \emph{Law of Quadratic Reciprocity}.
\begin{Theorem}{Law of Quadratic Reciprocity}{thm:quad-reciprocity}
  For any two distinct odd primes $p,q$ we have $$\Leg{p}{q}\cdot \Leg{q}{p}=(-1)^{\nicefrac{(p-1)(q-1)}{4}}$$
\end{Theorem}
\begin{proof}
  Let $\chi_1$ be the canonical additive character of $\bbF_p$ and $\eta$ is the quadratic character of $\bbF_p$. Then by \GSidentitieslem[(iv)] we gave $G^2(\eta,\chi_1)=(-1)^{\nicefrac{(p-1)}2}\cdot p$. So \begin{equation}\label{eq:quad-gauss-qpow}
    G^q(\eta,\chi_1)= G(\eta,\chi_1)\cdot \lt( G^2(\eta,\chi_1)\rt)^{\nicefrac{(q-1)}{2}}=(-1)^{\nicefrac{(p-1)(q-1)}{4}}\cdot p^{\nicefrac{(q-1)}2}\cdot G(\eta,\chi_1)
  \end{equation}Now we calculate $G^q(\eta,\chi_1)$ in different way. \begin{align*}
    G^q(\eta,\chi_1) & =\lt(\sum_{\alpha\in\bbF_p^*}\eta(\alpha)\cdot \chi_1(\alpha)\rt)^q\\ 
    & \equiv \sum_{\alpha\in\bbF_p^*}\eta^q(\alpha)\cdot \chi_q^q(\alpha)\bmod{q}\\ 
    & \equiv \sum_{\alpha\in\bbF_p^*}\eta(\alpha)\cdot \chi_q(\alpha)\bmod{q}\\ 
    & \equiv G(\eta,\chi_q)\bmod{q}\\ 
    & \equiv \eta(q)G(\eta,\chi_1)\bmod{q}\byl[(i)]{gaussian-sum-identity}\numberthis\label{eq:quad-gauss-qpow-modq}
  \end{align*}
  Now we already obtained $\Leg{q}{p}$ in the right hand side above and $(-1)^{\nicefrac{(p-1)(q-1)}{4}}$ in \eqref{eq:quad-gauss-qpow}. Since $\eta$ and $\chi_1$ are both non-trivial characters of $\bbF_p$ we have $|G(\eta,\chi_1)|^=p^{\nicefrac12}$  by \GSpropthm and so $G(\eta,\chi_1)\neq 0$. So now comparing \eqref{eq:quad-gauss-qpow-modq} with \eqref{eq:quad-gauss-qpow} we get \begin{align*}
    (-1)^{\nicefrac{(p-1)(q-1)}{4}}\cdot p^{\nicefrac{(q-1)}2}\cdot G(\eta,\chi_1) \equiv \eta(q)\cdot G(\eta,\chi_1)\bmod q   \iff  (-1)^{\nicefrac{(p-1)(q-1)}{4}}\cdot p^{\nicefrac{(q-1)}2}\equiv \eta(q)\bmod q
  \end{align*}
  By \FLTthm we have $p^{q-1}\equiv 1\bmod q$ as $p$ and $q$ are two distinct primes. So we have \begin{align*}
    & (-1)^{\nicefrac{(p-1)(q-1)}{4}}\cdot p^{\nicefrac{(q-1)}2}\equiv \eta(q)\bmod q\\ 
    \implies & (-1)^{\nicefrac{(p-1)(q-1)}{4}}\cdot p^{\nicefrac{(q-1)}2}\cdot p^{\nicefrac{(q-1)}2}\equiv \eta(q)\cdot p^{\nicefrac{(q-1)}2}\bmod q\\ 
    \implies & (-1)^{\nicefrac{(p-1)(q-1)}{4}}\equiv \eta(q)\cdot p^{\nicefrac{(q-1)}2}\bmod q\\ 
    \implies & (-1)^{\nicefrac{(p-1)(q-1)}{4}}\equiv \Leg{q}{p}\cdot \Leg{p}{q}\bmod q
  \end{align*}Since both sides are $\pm1$ we can remove the modulo $q$. And therefore we have the theorem.
\end{proof}

We close the section with an elementary but remarkably clean evaluation of the quadratic character summed against a quadratic polynomial. Its value depends only on the leading coefficient and the discriminant, and it will play a basic role in many later arguments -- most notably in extracting the quadratic character of the discriminant when counting solutions of quadratic equations over $\bbF_q$, and in the reformulation of Kloosterman sums in \autoref{sec:kloosterman}.
\begin{Theorem}{Quadratic Character Sum via Discriminant}{thm:eta-quadratic-discriminant}
  Let $f(X)=a_2X^2+a_1X+a_0\in\bbF_q[X]$ with $q$ odd and $a_2\neq 0$, let $\Dl=a_1^2-4a_0a_2$ denote the discriminant of $f$, and let $\eta$ be the quadratic character of $\bbF_q$ (extended by $\eta(0)=0$). Then $$\sum_{\alpha\in\bbF_q}\eta(f(\alpha))=\begin{cases}-\eta(a_2) & \text{if }\Dl\neq 0,\\ (q-1)\cdot \eta(a_2) & \text{if }\Dl=0.\end{cases}$$
\end{Theorem}
We will not prove this theorem now. We will discuss this again in \autoref{sec:character-sum-poly}. 